\section{Finite reproduction and route consequence}

The certificate evaluates four clocks with $L=5,7,11,13$.  In each case
$Q=2^jP_L$ is chosen so that $\log Q$ is already close to $L\log L$.  It locates the
first prime in each endpoint interval, verifies the two cutoff identities, enumerates
coprime multipliers, and checks
\[
 A_{4LQ}(L)=1,
 \qquad
 A_{4LQ}(2L-1)=1+\pi(2L-1)-\pi(L).
\]
The largest certified scale is
\[
 Q=257955735797760,\qquad
 p=257955735797867,\qquad
 r=505990097141761.
\]
The four exact raw mass ratios are
$2,3,4,4$.  These small values are finite reproduction only; divergence is supplied
by Theorem~\ref{thm:main}.

One implementation uses the deterministic unsigned-64-bit Miller--Rabin basis set and
constructs the primorial clock.  A separate implementation uses a different
deterministic basis theorem and rebuilds every record from the serialized scale data.
Normal and optimized Python modes agree.  An adversarial script removes one required
primorial factor from each clock and verifies that the low-row identity fails, while
the cutoff and universal-envelope boundary checks remain exact.  The record digest is
\begin{center}
\texttt{93fec725ce67b017dffcf4b04ecb4125dd67b241d6e5cd12f18015fb0ee0f15c}.
\end{center}

The route implication is narrow but decisive.  A collision-incidence estimate may be
converted to an incident-mass estimate only after row weights are controlled.  At
critical depth, raw support count does not supply that control.  The next minimal
question is therefore whether unit-row normalization yields a uniformly stable
collision operator, and only afterward whether such normalization is source-valid.
